% !TeX spellcheck = en_US
\documentclass[french]{yLectureNote}

\title{Algèbre linéaire 2}
\subtitle{Langage mathématique}
\author{Paulhenry Saux}
\date{\today}
\yLanguage{Français}

\professor{C.Dartyge}

\usepackage{graphicx}%----pour mettre des images
\usepackage[utf8]{inputenc}%---encodage
\usepackage{geometry}%---pour modifier les tailles et mettre a4paper
%\usepackage{awesomebox}%---pour les boites d'exercices, de pbq et de croquis ---d\'esactiv\'e pour les TP de PC
\usepackage{tikz}%---pour deiffner + d\'ependance de chemfig
\usepackage{tkz-tab}
\usepackage{chemfig}%---pour deiffner formules chimiques
\usepackage{chemformula}%---pour les formules chimiques en \'equation : \ch{...}
\usepackage{tabularx}%---pour dimensionner automatiquement les tableaux avec variable X
\usepackage{awesomebox}%---Pour les boites info, danger et autres
\usepackage{menukeys}%---Pour deiffner les touches de Calculatrice
\usepackage{fancyhdr}%---pour les en-t\^ete personnalis\'ees
\usepackage{blindtext}%---pour les liens
\usepackage{hyperref}%---pour les liens (\`a mettre en dernier)
\usepackage{caption}%---pour la francisation de la l\'egende table vers Tableau
\usepackage{pifont}
\usepackage{array}%---pour les tableaux
\usepackage{lipsum}
\usepackage{yFlatTable}
\usepackage{multicol}
\newcommand{\Lim}[1]{\lim\limits_{\substack{#1}}\:}
\renewcommand{\vec}{\overrightarrow}
\newcommand{\bz}{\overline{z}}
\DeclareMathOperator{\card}{card}
\renewcommand{\vec}{\overrightarrow}
\newcommand{\N}[0]{\mathbb{N}}
\newcommand{\R}[0]{\mathbb{R}}
\newcommand{\C}[0]{\mathbb{C}}
\newcommand{\dd}[0]{\mathrm{d}}
\begin{document}
\setcounter{chapter}{3}
\chapter{Diagonalisation}
On suppose E de dimension finie.
\section{Problème}
Quand on on change f de base dans E, on a la relation
\[A_{e2} = P_{e,e_2}A_{e}P_{e_2e}\]
\begin{definition}[Matrice diagonalisable]
f est diagonalisable s’il existe une base , telle que :
\[M(f)_{e_i} =
\begin{pmatrix}
a_{11}&0&\dots&0\\
0&a_{22}&\ddots&\vdots\\
\vdots&\ddots&\ddots&0\\
0&\dots&0&a_{nn}
\end{pmatrix}\]
\end{definition}
\begin{definition}[Matrice trigonalisable]
f est trigonalisable s’il existe une base , telle que :
\[M(f)_{e_i} =
\begin{pmatrix}
a_{11}&+&\dots&*\\
0&a_{22}&\ddots&\vdots\\
\vdots&\ddots&\ddots&*\\
0&\dots&0&a_{nn}
\end{pmatrix}\]
\end{definition}
\checkInfo{Comparaison avec la physique}{Ce problème est comparable à
celui de la mécanique  lorsqu’on essaie de trouver le repère dans lequel le problème
a la forme la plus simple et la plus pratique. Par exemple, si l’on étudie la chute
libre d’un corps il est naturel de prendre
un repère dont l’un des axes est vertical,
de manière à ce que le vecteur accélération
n’ait qu’une composante non nulle. Il serait contraire au bon sens de prendre les
axes obliques selon lesquels on doit décomposer le vecteur accélération.}
\section{Vecteur propre}
\begin{definition}[Vecteur propre]
Un vecteur \(v\in E\) est dit vecteur propre de f si \(\v \neq 0\) et \(\exists \lambda \in \R, f(v)=\lambda v\). On dit que \(\lambda\) est valeur propre de \(v\)
\end{definition}
\criticalInfo{Valeurs propres et vecteurs propres}{
\begin{itemize}
 \item Les valeurs propres peuvent \^etre nulles.
 \item Les vecteurs propres sont soit :
 \begin{itemize}
  \item des vecteurs du noyau
  \item des vecteurs qui ne changent pas de direction par f. Ainsi, la droite engendrée par un vecteur propre est invariante par f.
 \end{itemize}
\end{itemize}
}
\begin{theorem}[Vecteurs propres et diagonalisation]
 \(f \in L_E\) est diagonalisable si et seulement si il existe
une base de E formée de vecteurs propres.
\end{theorem}
\section{Recherche de vecteurs propres}
\begin{definition}[Polynome caractéristique]
Le polynome caractéristique de f est un polynome de degré n en \(\lambda\) : \(P_f(\lambda):= det(f-\lambda I_d)\)
\end{definition}
\begin{proposition}[Valeur propre de f et polynome caractéristique]
Les valeurs propres de f sont les racines du polynome caractéristique.
\end{proposition}
\begin{definition}[Spectre]
L’ensemble des valeurs propres de f est dit spectre de f et est noté \(Sp_K (f)\) avec K les réels ou les complexes.
\end{definition}
La précision du corps \(K\) est importante car parfois, on peut trouver des racines complexes.
\section{Recherche de vecteurs propres}
\warningInfo{Astuce}{La trace de la matrice dans la nouvelle base doit \^etre la m\^eme que celle de la matrice d'origine !}
\section{Caractérisation des endomorphismes diagonalisables}
\begin{definition}[Espace propre]
 \(E_{\lambda} = \{v\in E|f(v)=\lambda  v\}\) est l'espace propre correspondant à \(\lambda\)
\end{definition}
\begin{proposition}[Somme directe d'espace propre]
Soient \(\lambda_1,\dots,\lambda_n\) 2 à 2 disitncts. Alors les espaces propres \(E_{\lambda_1}, \dots, E_{\lambda_{n}}\) sont en somme directe.
\end{proposition}
\begin{theorem}[Propriétés équivalentes pour les espaces propres]

\end{theorem}

\end{document}


